\metatitle{Unicellular LLT polynomials and twin manifolds}
\metadescription{Unicellular LLT polynomials, their graph-coloring model,
Schur and elementary expansions, Frobenius characteristics, and the
twin-manifold dagger action.}


\section[unicellularLLT]{Unicellular LLT polynomials}

\begin{polydata}{unicellularLLT}
  Name   & Unicellular LLT polynomials \\
  Space    & Sym \\
  Basis    & False \\
  Rating   & 4 \\
  Bib      & HaglundHaimanLoehr2005\\
  Year     & 2005\\
  Symbol   & $\LLT_\avec(\xvec;q)$ \\
  Keywords & fillings, schur-positive, murnaghan-nakayama \\
  Category & Schur \\
  PositiveIn & schurS | Blasiak2016 | scope=max_i a_i <= 2; proof=explicit combinatorial formula \\
  PositiveIn & kSchur | Miller2019 | scope=bandwidth at most 3; note=3-Schur positivity \\
  PositiveIn & elementaryE | AlexanderssonSulzgruber2020x | scope=unit interval graphs after q -> q+1 \\
\end{polydata}


The \defin{unicellular LLT polynomials} are the one-cell-per-component
subfamily of the \hyperref[LLT]{LLT polynomials}.  In the tuple model, this
means that the indexing tuple consists of single boxes.  A more convenient
description, observed in \cite{CarlssonMellit2017,AlexanderssonPanova2018},
uses \hyperref[chromaticQuasisymmetricUnitIntervalGraph]{unit interval graphs}
and their area sequences.


\subsection[unicellularLLTGraphModel]{The graph-coloring model}

Let $\avec=(a_1,\dotsc,a_n)$ be an area sequence of a Dyck path, and let
$\Gamma_\avec$ be the corresponding natural unit interval graph on vertex set
$[n]$.  We use the same graph and ascent statistic as on the
\hyperref[chromaticQuasisymmetricDefinition]{chromatic quasisymmetric function
page}.  The vertices are ordered as
\[
1\lt 2\lt \dotsb \lt n.
\]
For a coloring $ \kappa:[n]\to \setP$, 
an edge $\{i,j\}$ with $i\lt j$ is an \defin{ascent} of $\kappa$ if
$ \kappa(i)\lt\kappa(j)$.
We write $\asc(\kappa)$ for the number of such edges.  The \defin{unicellular LLT
polynomial} is
\[
\LLT_\avec(\xvec;q)
\coloneqq
\sum_{\kappa:[n]\to\setP}
x_{\kappa(1)}\dotsm x_{\kappa(n)}q^{\asc(\kappa)}.
\]
All colorings are allowed. This is the main difference from the corresponding
\hyperref[chromaticQuasisymmetricDefinition]{chromatic quasisymmetric function}.

The graph-coloring model is equivalent to the one-cell tuple definition of LLT
polynomials in the Haglund--Haiman--Loehr model
\cite{HaglundHaimanLoehr2005}.  The explicit bijection between tuples of
single boxes and Dyck diagrams is described in \cite{AlexanderssonPanova2018}.

\begin{example*}[Bijection between tuples and Dyck diagrams]
The $3$-tuple of skew shapes $3211/211$, $4221/321$, $421/31$ gives a
unicellular LLT polynomial.  The boxes are labeled in reading order, and the
corresponding boxes are labeled in the Dyck diagram with area sequence
$011101012$.
\begin{figure}
\begin{ytableau}
&  &  &  &  &   &  &  & *(lightblue) 3 &   &   &  \\
&  &  &  &  &   &  &  & *(lightblue) & *(lightblue) 5 &   &  \\
&  &  &  &  &   &  &  & *(lightblue) & *(lightblue) & *(lightblue) & *(lightblue) 9\\
&  &  &  & *(lightblue) 1 &   &  &  &  &  &  & \\
&  &  &  & *(lightblue) & *(lightblue) 4 &  &  &  &  &  & \\
&  &  &  & *(lightblue) & *(lightblue)  &  &  &  &  &  & \\
&  &  &  & *(lightblue) & *(lightblue)  &*(lightblue)  & *(lightblue) 8 &  &  &  & \\
*(lightblue)2&  &  &  &  &  &  &  &  &  &  & \\
*(lightblue) &  &  &  &  &  &  &  &  &  &  & \\
*(lightblue) & *(lightblue) 6 &  &  &  &  &  &  &  &  &  & \\
*(lightblue) & *(lightblue) & *(lightblue) 7 &  &  &  &  &  &  &  &  &
\end{ytableau}
\begin{ytableau}
*(lightgray) & *(lightgray) & *(lightgray) & *(lightgray) & *(lightgray) & *(lightgray) & & & *(yellow) 9 \\
*(lightgray) & *(lightgray) & *(lightgray) & *(lightgray) & *(lightgray) & *(lightgray) &  & *(yellow) 8 \\
*(lightgray) & *(lightgray) & *(lightgray) & *(lightgray) & *(lightgray) & *(lightgray) & *(yellow) 7 \\
*(lightgray) & *(lightgray) & *(lightgray) & *(lightgray) & & *(yellow) 6 \\
*(lightgray) & *(lightgray) & *(lightgray) & *(lightgray) & *(yellow) 5 \\
*(lightgray) & *(lightgray) &  & *(yellow) 4 \\
*(lightgray) &  & *(yellow) 3 \\
 & *(yellow) 2 \\
*(yellow) 1
\end{ytableau}
\end{figure}
In general, if the tuple consists of possibly disconnected ribbons of size at
most $k$, then the entries in the corresponding area sequence are less than
$k$.
\end{example*}


\subsection[unicellularLLTChromaticRelation]{Relation with chromatic quasisymmetric functions}

The same graph and ascent statistic control both families.  The chromatic
quasisymmetric function imposes properness; the LLT polynomial does not.  In
\cite{CarlssonMellit2017}, \name{Erik Carlsson} and \name{Anton Mellit} prove
the plethystic relation
\[
(q-1)^{-n}\LLT_\avec[\xvec(q-1);q]
=
\chrom_{\Gamma_\avec}(\xvec;q).
\]

\subsection[unicellularLLTProperties]{Basic properties}

The complete graphs on $n$ vertices are unit interval graphs.  For the complete
graph $K_n$, one has
\[
\LLT_{K_n}(\xvec;q)
=
\macdonaldH_{(n)}(\xvec;q,t)
=
\sum_{T\in\SYT(n)}q^{\cocharge(T)}\schurS_{\lambda(T)}(\xvec).
\]
Here $\macdonaldH_{(n)}(\xvec;q,t)$ is a
\hyperref[macdonaldH]{modified Macdonald polynomial}, and the last sum is a
\hyperref[hallLittlewoodH]{modified Hall--Littlewood polynomial}.

In \cite{AlexanderssonPanova2018}, it is proved that for any area sequence
$\avec$ of length $n$,
\[
\omega \LLT_{\avec}(\xvec;q)
=
q^{a_1+\dotsb+a_n}\LLT_{\avec^T}(\xvec;1/q),
\]
where $\avec^T$ denotes the transpose of the Dyck diagram.  From the definition,
\[
\LLT_{\avec^T}(x_1,x_2,\dotsc,x_n;q)
=
\LLT_{\avec}(x_n,x_{n-1},\dotsc,x_1;q),
\]
and since the polynomials are symmetric, it follows that
$\LLT_{\avec^T}(\xvec;q)=\LLT_{\avec}(\xvec;q)$.

The unicellular LLT polynomials are also evaluations of certain traces of
\hyperref[heckeAlgebra]{Hecke algebras}; see \cite{MoralesSkanderaWang2024}.
There is a close connection with Kazhdan--Lusztig polynomials.


\subsection[unicellularLLTRecursions]{Recursions}

There are several recursions and linear relations related to unicellular LLT
polynomials.  For example, \name{Seung Jin Lee} gives a three-term recursion in
\cite{Lee2020}.  Short bijective proofs of these recursions are given in
\cite{Alexandersson2019llt}.

Lee's recursions are generalized further by \name{Christopher Miller}; see
\cite[Eq. (2.2)]{Miller2019}.  For more generalizations and results on linear
relations between LLT polynomials, see \cite{Tom2020x} and \cite{Keating2021x}.


\subsection[unicellularLLTSchurExpansion]{Schur expansion}

It is a major open problem to find a combinatorial rule for the Schur expansion
of $\LLT_\avec(\xvec;q)$.  For so-called melting lollipop graphs, a
combinatorial formula for the Schur expansion is given in \cite{HuhNamYoo2020}.
\name{Victor Wang} extends this direction to two-headed melting lollipop
graphs by giving a formula expanding the corresponding unicellular LLT
polynomials into ribbon Schur functions \cite{WangLollipopLLT2025x}.

\begin{example*}[Schur expansions for length 3 area sequences]
For area sequences of length $3$, one has
\begin{align*}
\LLT_{000}(\xvec;q)&=\schurS_{3}+2\schurS_{21}+\schurS_{111}, \\
\LLT_{001}(\xvec;q)&=\schurS_{3}+(1+q)\schurS_{21}+q\schurS_{111}, \\
\LLT_{011}(\xvec;q)&=\schurS_{3}+2q\schurS_{21}+q^2\schurS_{111}, \\
\LLT_{012}(\xvec;q)&=\schurS_{3}+(q+q^2)\schurS_{21}+q^3\schurS_{111}.
\end{align*}
\end{example*}

Since LLT polynomials are $q$-deformations of
\hyperref[schurLittlewoodRichardson]{Littlewood--Richardson coefficients}
\cite{LeclercThibon2000}, the unicellular case suggests the following
standard-tableau formulation.

\begin{conjecture}[See \cite{LeclercThibon2000}]
There is a statistic $c_\avec(T)$ on
\hyperref[prelimTableaux]{standard Young tableaux}, depending on $\avec$, such
that
\[
\LLT_{\avec}(\xvec;q)
=
\sum_{T\in\SYT(n)}q^{c_\avec(T)}\schurS_{\lambda(T)}.
\]
By comparing the coefficients of $x_1x_2\dotsm x_n$ on both sides, one obtains
\[
\sum_{\sigma\in\symS_n} q^{\asc_\avec(\sigma)}
=
\sum_{\lambda\vdash n}f^\lambda
\sum_{T\in\SYT(\lambda)}q^{c_\avec(T)},
\]
where the left-hand side sums over all vertex colorings of $\Gamma_\avec$
using the $n$ distinct colors $1,\dotsc,n$, and $f^\lambda$ is the number of
\hyperref[prelimTableaux]{standard Young tableaux} of shape $\lambda$.
\end{conjecture}

\begin{theorem}[See \name{Jonah Blasiak} \cite{Blasiak2016}]
Whenever $\nuvec=(\nu^1,\nu^2,\nu^3)$ is a $3$-tuple of possibly disconnected
skew shapes, $\LLT_{\nuvec}(\xvec;q)$ is Schur-positive with an explicit
combinatorial formula.  In the unicellular case, this includes the area
sequences indexing $\LLT_{\avec}(\xvec;q)$ with $\max_i a_i\leq 2$.
\end{theorem}

In his PhD thesis, \name{Christopher Miller} shows that LLT polynomials with
bandwidth at most $3$ are \hyperref[kSchur]{$3$-Schur-positive}, which is
stronger than Schur positivity \cite{Miller2019}.  The bandwidth-at-most-two
case was treated earlier by \name[S.J. Lee]{Seung Jin Lee} \cite{Lee2020}.

In \cite{ChoHuh2019}, a combinatorial formula is given for the hook
coefficients, that is, the coefficients of $\schurS_{\lambda}$ where
$\lambda$ has the form $(m,1,1,\dotsc,1)$.


\subsection[unicellularLLTPExpansion]{Power-sum expansion}

From the \hyperref[chromaticQuasisymmetricPowerSumExpansion]{power-sum
expansion of chromatic quasisymmetric functions}, one can derive a formula for
$\LLT_\avec(\xvec;q)$ in the power-sum basis.  Alternatively, the general
machinery in \cite{AlexanderssonSulzgruber2019} gives the following shifted
power-sum expansion; see also the related result on
\hyperref[pPartitionPsiExp]{P-partitions}.

Given a poset $P$ on $n$ elements, let $\opsurj(P)$ be the set of
order-preserving surjections $f:P\to[k]$ for some $k$.  The \defin{type} of
such a surjection is
\[
\alpha(f)\coloneqq (|f^{-1}(1)|,|f^{-1}(2)|,\dotsc,|f^{-1}(k)|).
\]
Let $\opsurj_\alpha(P)\subseteq\opsurj(P)$ be the set of surjections of type
$\alpha$, and let $\opsurj^\ast_\alpha(P)\subseteq\opsurj_\alpha(P)$ be the set
of surjections $f\in\opsurj_\alpha(P)$ such that, for each $j\in[k]$,
$f^{-1}(j)$ is an induced subposet of $P$ with a unique minimal element.

\begin{theorem}[Alexandersson--Sulzgruber \cite{AlexanderssonSulzgruber2019}]
Let $O(\avec)$ denote the set of orientations of the unit interval graph given
by $\avec$.  Then
\[
\omega \LLT_{\avec}(\xvec;q+1)
=
\sum_{\theta\in O(\avec)}q^{\asc(\theta)}
\sum_{\lambda\vdash n}
\frac{\powerSum_\lambda(\xvec)}{z_\lambda}
|\opsurj^\ast_\lambda(P(\theta))|,
\]
where $P(\theta)$ is the poset on $[n]$ given by the transitive closure of the
ascending edges in $\theta$.
\end{theorem}

\begin{example*}[Power-sum expansion for $\avec=0121$ with orientations]
The area sequence $\avec=(0,1,2,1)$ gives the expansion
\begin{align*}
\omega \LLT_{0121}(\xvec;q+1)
= {}&(2q^3+q^4)\frac{\powerSum_4}{z_4}
+(4q^2+4q^3+q^4)\frac{\powerSum_{31}}{z_{31}} \\
&+(2q^2+2q^3+q^4)\frac{\powerSum_{22}}{z_{22}}
+(8q+12q^2+6q^3+q^4)\frac{\powerSum_{211}}{z_{211}} \\
&+(24+48q+34q^2+10q^3+q^4)
\frac{\powerSum_{1111}}{z_{1111}}.
\end{align*}

The following two orientations $\theta_1,\theta_2$ contribute to
$2q^3\powerSum_{22}/z_{22}$.  By convention, only the ascending edges are
shown.
\begin{figure}
\begin{ytableau}
*(lightgray) & *(lightgray) & \to & *(yellow) 4 \\
 & \to & *(yellow) 3 \\
\to & *(yellow) 2 \\
*(yellow) 1
\end{ytableau}
\begin{ytableau}
*(lightgray) & *(lightgray) & \to & *(yellow) 4 \\
\to &  & *(yellow) 3 \\
\to & *(yellow) 2 \\
*(yellow) 1
\end{ytableau}
\end{figure}
For each orientation, the order-preserving surjection $f$ defined by
$f(1)=f(2)=1$ and $f(3)=f(4)=2$ is the only one satisfying the conditions for
membership in $\opsurj^\ast_{22}(P(\theta))$.  The minimal elements of
$f^{-1}(1)$ and $f^{-1}(2)$ are $1$ and $3$, respectively.

For $\theta_2$, the map $g(1)=g(3)=1$ and $g(2)=g(4)=2$ is also an
order-preserving surjection of type $(2,2)$, but the induced subposet
$g^{-1}(2)\subset P(\theta_2)$ has both elements as minimal elements.  Hence
$g$ is not an element of $\opsurj^\ast_{22}(P(\theta_2))$.
\end{example*}

A similar-looking formula for
\hyperref[chromaticQuasisymmetricPowerSumExpansion]{chromatic symmetric
functions} is given in \cite[Prop. 5.2]{BernardiNadeau2020}.


\subsection[unicellularLLTEExpansion]{Elementary symmetric expansion}

Together with \name{Greta Panova}, we conjectured in
\cite{AlexanderssonPanova2018} that for any unit interval graph $\avec$, the
polynomial $\LLT_\avec(\xvec;q+1)$ is positive in the elementary symmetric
function basis.  The formula below was later proved by \name{Per Alexandersson} and
\name{Robin Sulzgruber}.

Let $\theta\in O(\avec)$, and for each $i\in[n]$ define the
\defin{highest reachable vertex} by
\[
\mathrm{hrv}(i)\coloneqq \max\{j\in[n]: i\leq_{P(\theta)}j\}.
\]
Here $P(\theta)$ is the poset generated by the ascending edges of
$\theta$.  Note that $\mathrm{hrv}(i)\geq i$ for all $i\in[n]$.  Finally let
\[
\pi(\theta)=(|\mathrm{hrv}^{-1}(1)|,\dotsc,|\mathrm{hrv}^{-1}(n)|),
\]
with zero parts omitted when it is used as an index of $\elementaryE$.

\begin{theorem}[Alexandersson--Sulzgruber \cite{AlexanderssonSulzgruber2020x}]
For any unit interval graph $\avec$ with $n$ vertices,
\[
\LLT_\avec(\xvec;q+1)
=
\sum_{\theta\in O(\avec)}q^{\asc(\theta)}
\elementaryE_{\pi(\theta)}(\xvec).
\]
\end{theorem}

This is an analogue of the
\hyperref[chromaticQuasisymmetricElementaryExpansion]{Stanley--Stembridge
conjecture} for \hyperref[chromaticQuasisymmetricDefinition]{chromatic
quasisymmetric functions}.  A similar formula holds for vertical-strip LLT
polynomials.  The path graph case was proved in
\cite{AlexanderssonPanova2018}, and the complete graph case can be proved by a
recursion.  For melting lollipop graphs, a similar-looking
$\elementaryE$-expansion is proved in \cite{Alexandersson2019llt}.

\begin{example}
In the diagram with $\avec=012223$, the orientation $\theta\in O(\avec)$ is
shown by marking the ascending edges with $\to$.
\begin{figure}
\begin{ytableau}
*(lightgray)   & *(lightgray)   &    &      & \to   & *(yellow) 6 \\
*(lightgray)   & *(lightgray)   &     &     & *(yellow) 5\\
*(lightgray)  &    &  \to   & *(yellow) 4\\
  \to  &     & *(yellow) 3\\
  \to  & *(yellow) 2\\
*(yellow) 1
\end{ytableau}
\end{figure}
The highest reachable vertices for $1,2,\dotsc,6$ are
$4,2,4,4,6,6$, respectively.  Hence
$\pi(\theta)=(0,0,1,0,3,0,2)$, and this orientation contributes
$q^5\elementaryE_{321}(\xvec)$ to $\LLT_{012223}(\xvec;q+1)$.

The full expansion is
\begin{align*}
\LLT_{012223}(\xvec;q+1)
= {}&(24q^5+60q^6+62q^7+33q^8+9q^9+q^{10})\elementaryE_6 \\
&+(8q^4+12q^5+6q^6+q^7)\elementaryE_{33} \\
&+(20q^4+36q^5+25q^6+8q^7+q^8)\elementaryE_{42} \\
&+(40q^4+96q^5+94q^6+46q^7+11q^8+q^9)\elementaryE_{51} \\
&+(4q^3+4q^4+q^5)\elementaryE_{222} \\
&+(36q^3+50q^4+24q^5+4q^6)\elementaryE_{321} \\
&+(34q^3+69q^4+56q^5+21q^6+3q^7)\elementaryE_{411} \\
&+(19q^2+17q^3+4q^4)\elementaryE_{2211} \\
&+(20q^2+28q^3+15q^4+3q^5)\elementaryE_{3111} \\
&+(10q+6q^2+q^3)\elementaryE_{21111}
+\elementaryE_{111111}.
\end{align*}
\end{example}

A refinement to LLT cumulants, where certain connectivity requirements are
imposed, is given in \cite{Kowalski2020x}; see also
\cite{DolegaKowalski2021x}.  Another refinement is given by
\name{Alex Abreu} and \name{Antonio Nigro} in \cite{AbreuNigro2021}, where
they consider symmetric functions defined as sums over increasing forests.
The unicellular LLT polynomials can be expressed as sums over these forests.


\section[unicellularLLTRepresentationTheory]{Representation theory}

The polynomial $\LLT_\avec(\xvec;q)$  is the Frobenius characteristic of
\emph{twin cohomology} with \emph{the dagger action}. 

\subsection[unicellularLLTFrobeniusTarget]{Frobenius characteristic}

The polynomial $\LLT_\avec(\xvec;q)$ is symmetric and Schur-positive. 
Hence it can be interpreted as a graded Frobenius characteristic:
\[
\LLT_\avec(\xvec;q)
=
\sum_{d\geq 0}q^d\frobChar(M_\avec^d)
\]
for some graded $\symS_n$-module
\[
M_\avec=\bigoplus_{d\geq 0}M_\avec^d.
\]
The question is to construct a natural module $M_\avec$ with this graded
Frobenius characteristic.

\name{Mathieu Guay-Paquet} used Hopf-algebra methods to connect the
Shareshian--Wachs theorem with unicellular LLT polynomials
\cite{GuayPaquet2016x}.  The twin-manifold story gives a cohomological
realization:
\[
\sum_{d\geq 0} q^d
\frobChar\left(H^{2d}(\mathrm{Twin}_\avec)\right)
=
\LLT_\avec(\xvec;q).
\]
This is the LLT analogue of the Hessenberg dot-action theorem.
\name{Young-Hoon Kiem} and \name{Donggun Lee} give a direct geometric proof
of this statement by comparing twin manifolds through a roof manifold and
showing that their cohomology satisfies the modular law for unicellular LLT
polynomials \cite[Thm.~5.4]{KiemLee2024Twin}.  This proof does not pass
through the Shareshian--Wachs conjecture for chromatic quasisymmetric
functions.
The Hilbert series of this graded module is obtained by pairing with
$\completeH_{1^n}$:
\[
\left\langle \LLT_\avec(\xvec;q),\completeH_{1^n}\right\rangle
=
\sum_{\sigma\in\symS_n}q^{\asc(\sigma)}.
\]
The right-hand side is the ascent enumerator over colorings using all $n$
colors exactly once, equivalently over permutations.

There is also a complementary finite-group interpretation.  In
\cite[Thm.~5.1]{Gagnon2024}, \name{Lucas Gagnon} realizes vertical-strip LLT
polynomials by inducing characters from the unipotent upper triangular group
to $\GL_n(\mathbb{F}_q)$ and then recording the unipotent constituents.


\subsection[unicellularLLTTwinGKM]{The twin GKM model}

The twin-manifold GKM model is close to the regular semisimple Hessenberg GKM
model, with two differences.

First, the congruence labels are fixed \hyperref[rootSystem]{roots}.
The fixed points are indexed by permutations $v\in\symS_n$, and for each edge $\{i,j\}$ of
$\Gamma_\avec$ we have a moment-graph edge
\[
v\longleftrightarrow v s_{ij}.
\]
However, the GKM divisibility relation is
\[
p_v-p_{v s_{ij}}\in (t_i-t_j),
\]
not
\[
p_v-p_{v s_{ij}}\in (t_{v(i)}-t_{v(j)}).
\]
Thus the edge label depends on the graph edge $\{i,j\}$ itself, rather than
on the image of the edge under the fixed point $v$.

Second, the group action is the \defin{dagger action}.  
It acts by moving the fixed-point labels, without permuting the polynomial variables:
\[
(\sigma\dagger p)_v
\coloneqq
p_{\sigma^{-1}v}.
\]
This should be contrasted with Tymoczko's dot action, where $\sigma$ also acts
on the variables $t_1,\dotsc,t_n$.

The equivariant twin GKM module is therefore
\[
H_T^\ast(\mathrm{Twin}_\avec)
=
\left\{
(p_v)_{v\in\symS_n}
:
p_v-p_{v s_{ij}}\in (t_i-t_j)
\text{ for all } v \text{ and } \{i,j\}\in E(\Gamma_\avec)
\right\}.
\]
Ordinary cohomology is obtained by setting $t_1,\dotsc,t_n$ to zero:
\[
H^\ast(\mathrm{Twin}_\avec)
\cong
H_T^\ast(\mathrm{Twin}_\avec)
\otimes_{\setC[t_1,\dotsc,t_n]}\setC.
\]
The dagger action descends to ordinary cohomology.
Taking the \hyperref[frobeniusCharacteristic]{graded Frobenius characteristic}
then gives the unicellular LLT polynomial.


\section[unicellularLLTNonCom]{Non-commutative unicellular LLT polynomials}

\begin{polydata}{unicellularLLTNonCom}
  Name   & Non-commutative unicellular LLT polynomials \\
  Space    & WQSym \\
  Basis    & False \\
  Rating   & 1 \\
  Bib      & NovelliThibon2019 \\
  Year     & 2019\\
  Symbol   & $\LLTnc_\nuvec(\xvec;q)$ \\
  Keywords & gessel-nc-positive, murnaghan-nakayama \\
  Category & Other \\
\end{polydata}


\name[J.-C. Novelli]{Jean-Christophe Novelli} and
\name[J.-Y. Thibon]{Jean-Yves Thibon} introduce a non-commutative lift of the
unicellular LLT polynomials in \cite[Thm. 6.2]{NovelliThibon2019}.  They show
that these are positive in a non-commutative analogue of the
\hyperref[gessel]{Gessel quasisymmetric basis}, and that they satisfy a
non-commutative \hyperref[plethysm]{plethystic} relation with
\hyperref[chromaticQuasisymmetricLLT]{chromatic quasisymmetric functions}.
